\documentclass{article}
\input{../../lib.sty}
\allowdisplaybreaks

\title{\texttt{impl TopKMeasure for MultiDP}}
\author{\MichaelShoemateAuthor}
\begin{document}
\maketitle

This document proves soundness of \rustdoc{measurements/noisy_top_k/trait}{TopKMeasure} for \texttt{MultiDP}.

\section{Hoare Triple}
\subsection*{Precondition}
\subsubsection*{Compiler-verified}
\begin{itemize}
    \item Method \texttt{privacy\_map}: types are consistent with the PureDP and zCDP top-$k$ accounting implementations.
\end{itemize}

\subsubsection*{Caller-verified}
\begin{itemize}
    \item Method \texttt{privacy\_map}: \texttt{d\_in} is non-null and non-negative, and \texttt{scale} is non-null and non-negative.
\end{itemize}

\subsection*{Postcondition}
\begin{theorem}
    The implementation is consistent with the associated items in the \rustdoc{measurements/noisy\_top\_k/trait}{TopKMeasure} trait.
    For any $x, x'$ where $d_\mathrm{in} \ge d_\mathrm{Range}(x, x')$, \texttt{privacy\_map} returns a \texttt{PrivacyGuarantee} containing a valid bound for the selected top-$k$ mechanism.
\end{theorem}

\begin{proof}
    The accounting preference selects the mechanism in both sampling and privacy accounting.
    An automatic preference is rejected because the score metric does not distinguish the algorithms.
    At zero sensitivity the mechanism is identical on adjacent inputs; at positive sensitivity and zero scale a vacuous bound is valid.
    Otherwise, for permute-and-flip, the implementation uses the established PureDP bound
    $\epsilon = k d_\mathrm{in}/\mathrm{scale}$ and embeds it as a pure-DP certificate.
    For the exponential mechanism (equivalent to Gumbel noisy selection), it uses the established zCDP bound
    $\rho = k(d_\mathrm{in}/\mathrm{scale})^2/8$.
    It also satisfies the same pure-DP bound as permute-and-flip:
    in each peeling round the exponential mechanism's privacy loss has range at most
    $d_\mathrm{in}/\mathrm{scale}$, since the log ratio of normalizers is
    independent of the selected index and the range metric bounds the score differences.
    The interval containing the privacy loss must contain zero, since both
    output distributions sum to one. Thus its absolute privacy loss is bounded by
    $d_\mathrm{in}/\mathrm{scale}$.
    For the RDP certificate, write $\nu=d_\mathrm{in}/\mathrm{scale}$ and let
    $\alpha>1$. A likelihood ratio with log-range at most $\nu$ has mean one under
    the denominator distribution. By convexity of $r\mapsto r^\alpha$, its
    $\alpha$-moment is upper-bounded by that of a two-point distribution on the
    endpoints of an interval of likelihood ratios with ratio $e^\nu$ and mean one.
    If $q$ is its mass on the upper endpoint, this moment equals
    \[
        \frac{1+qB}{(1+qA)^\alpha},\qquad
        A=e^\nu-1,\quad B=e^{\alpha\nu}-1.
    \]
    For $\nu>0$ the derivative of its logarithm has the sign of
    $B-\alpha A-(\alpha-1)ABq$. Maximizing over $q$ gives the upper bound
    \[
        D_\alpha \le \frac{\log(B/(\alpha A))}{\alpha-1}
                    -\log\frac{\alpha(B-A)}{(\alpha-1)B}.
    \]
    Factoring the exponentials out of $A$, $B$, and $B-A$ yields the log-space
    expression in the implementation. Interval arithmetic encloses this bound;
    if its evaluation fails, the pure-DP bound $\nu$ is valid at every order.
    At $\nu=0$ the RDP bound is zero.

    The bounds compose over the $k$ peeling rounds. All three certificates are
    attached for the exponential mechanism when the pure-DP bound is finite,
    and the pure-DP certificate for permute-and-flip when finite.
    Directed arithmetic rounds the final multiplication upward.
\end{proof}

\bibliographystyle{plain}
\bibliography{references.bib}

\end{document}
